% /root/monodromy, 19 November 2026; complete conditional reduction independently audited by /root.
\subsection{A single geometric gate for the full product groups}
\label{subsec:product-irreducible-hasse-reduction}
\status{Reducedness and the complete reduction proved and independently audited by the parent. The geometric hypothesis is discharged for every degree at least four in Theorem~\ref{thm:product-Hasse-irred}; the resulting unconditional all-degree images are Corollary~\ref{thm:product-char2-all-full-groups}.}

\begin{proposition}\label{prop:product-hasse-reduced}
For the fixed surface $E\times E$ and embedding considered above,
the full Hasse polynomial $h_d(s)=\det M_d(s)$ is geometrically
reduced for every $d\ge1$.
\end{proposition}

\begin{proof}
At the geometric generic point of any irreducible Hasse component,
Theorem~\ref{thm:product-all-degree-transvections} gives stable rank
$b-1$ for the full coherent matrix, where $b=9d^2$.
Thus $M$ and $M^{[2]}M$ both have rank $b-1$.
Let $u$ span $\ker M$ and let $\ell$ span the left kernel.
Since $\ker M^{[2]}$ is spanned by $u^{[2]}$, equality of the two
ranks says that $u^{[2]}$ is not in $\im M$, or equivalently
$\ell(u^{[2]})\ne0$.

The symmetric Cartier multiplication pairing gives
$M(s)w=M(w)s$ and $M(w)w=w^{[2]}$.
The adjugate at this rank-$b-1$ point is a nonzero scalar multiple
of $u\ell$.  Differentiating the determinant in the coefficient
direction $u$ therefore gives a nonzero scalar multiple of
\[
                     \ell M(u)u=\ell(u^{[2]})\ne0.
\]
Each component is consequently generically smooth and occurs
with multiplicity one.  The repeated factors in some retained
pencil determinants reflect tangency of those pencils, not a
nonreduced global Hasse divisor.
\end{proof}

\begin{lemma}[Frobenius-periodic reflexive sheaves]
\label{lem:periodic-reflexive-projective-trivial}
Let $P$ be projective space over an algebraically closed field of
characteristic two.  If a coherent reflexive sheaf $\mathcal S$
has an isomorphism $F^*\mathcal S\simeq\mathcal S$, then it is a
trivial vector bundle with its constant unit-Frobenius structure.
\end{lemma}

\begin{proof}
Fitting ideals commute with base change.  In a local ring of $P$,
each Fitting ideal $I$ therefore satisfies $I=I^{[2]}$.
If $I$ is proper, this gives $I\subseteq\mathfrak m I$, so
Nakayama makes $I=0$.  The usual Fitting-ideal criterion, or a
minimal presentation, now shows that the sheaf is locally free
of its generic rank.  Its Frobenius isomorphism gives a finite
etale Lang frame cover.  Projective space has trivial geometric
fundamental group, as proved in
Lemma~\ref{lem:char2-defect-one-inertia}, so this frame cover has
a section.  The resulting frame is Frobenius-fixed.
\end{proof}

\begin{lemma}[Determinant allocation]
\label{lem:irreducible-hasse-quotients-trivial}
Let $P=\mathbf P^m$ over an algebraically closed field of
characteristic two, and let $\mathcal E=\Omega_P^1$ carry a
generically invertible Frobenius-linear map.  Suppose its Hasse
divisor is reduced and irreducible.
Every nonzero proper invariant subspace of its Lang representation
has trivial quotient representation.
\end{lemma}

\begin{proof}
Let $U$ be the invertibility open and extend the unit-Frobenius
subbundle of an invariant subspace by saturation in $\mathcal E$.
This gives a coherent exact sequence
\[
 0\longrightarrow\mathcal S\longrightarrow\mathcal E
       \longrightarrow\mathcal Q\longrightarrow0,
\]
where $\mathcal Q$ is torsion-free and $\mathcal S$ is reflexive.
The latter assertion follows locally from the depth lemma:
the kernel of a vector bundle mapping to a torsion-free sheaf
on a regular scheme satisfies the depth-two condition.
Frobenius preserves the saturated subsheaf, because it does so
at the generic point.  Flatness of Frobenius on the regular
scheme preserves reflexivity (regular sequences may equivalently
be replaced by their squares).  It therefore induces regular
maps on $\mathcal S$ and, after reflexive extension, on
$\mathcal Q^{**}$.

In codimension one all terms are vector bundles, and determinants
of these maps give effective divisors satisfying
\[
                     D_{\mathcal E}=D_{\mathcal S}+D_{\mathcal Q}.
\]
If $D_{\mathcal S}=0$, the map
$F^*\mathcal S\to\mathcal S$ is an isomorphism in codimension
one.  Both sheaves are reflexive, so it is an isomorphism
everywhere.  Lemma~\ref{lem:periodic-reflexive-projective-trivial}
makes $\mathcal S$ a nonzero trivial bundle.  Its inclusion in
$\Omega_P^1$ contradicts $H^0(P,\Omega_P^1)=0$, which follows
at once from the Euler sequence.

The reduced irreducible divisor must consequently be allocated
entirely to $\mathcal S$, and $D_{\mathcal Q}=0$.
The same argument makes $\mathcal Q^{**}$ a trivial unit-Frobenius
bundle.  On $U$ it is the original quotient bundle, because
that quotient is already locally free there.  The composite
$\mathcal E\to\mathcal Q\to\mathcal Q^{**}$ is Frobenius-equivariant
and exhibits the quotient local system as constant.
\end{proof}

\begin{lemma}[A fixed natural block in every higher degree]
\label{lem:product-degree-three-block}
For every $d\ge3$, the geometric group on the fixed surface
contains a subgroup acting as the natural
$\GL_{80}(\mathbf F_2)$ on an $80$-dimensional direct summand
and trivially on a complementary summand.
\end{lemma}

\begin{proof}
The assertion for $d=3$ has already been proved.  For $d>3$ fix
$d-3$ ordinary smooth members $D_i\in|L|$ in general transverse
position, with no triple points, and put $D=\bigcup D_i$.
This fixed connected nodal curve is ordinary.
Vary an ordinary degree-three curve $C$, transverse to every
$D_i$ and avoiding all nodes of $D$.
The resulting curves $C\cup D$ are ordinary by the union sequence,
so their Lang sheaf is a restriction of the full ordinary
coherent family in degree $d$.  Its image is a subgroup of the
degree-$d$ geometric group.  The allowed $C$ form a dense open
of the degree-three ordinary parameter space; hence their
relative degree-three image remains $\GL_{80}(\mathbf F_2)$.

The union sequence and the constant ambient
$H^1(X,\mathbf F_2)$ give a filtration on this boundary
representation $V$,
\[
 0\subset T\subset T_0\subset V,
 \qquad T_0/T\text{ constant},\qquad V/T_0=Q,
 \quad\dim Q=80.
\]
Here $T$ is the graph term from $C\cap D$ and $Q$ is the relative
degree-three local system.  More explicitly, the quotient by
$T$ is
$(H^1(C,\mathbf F_2)\oplus H^1(D,\mathbf F_2))/H^1(X,\mathbf F_2)$;
the inverse image of the constant $H^1(D,\mathbf F_2)$ is $T_0$.
The ambient maps are injective, so these identifications are exact.

Each $C\cap D_i$ consists of $3L^2=54$ points.
Thus the image on $T$ is a subgroup of a product of $S_{54}$'s.
The image on $T_0$ is an extension of such a subgroup by a
subgroup of the additive two-group
$\Hom(T_0/T,T)$.  Every prime in its order is at most $53$.
In contrast $\GL_{80}(\mathbf F_2)$ is simple and its order is
divisible by $4278255361$.

We recall an elementary simplicity proof to specify this input.
The group $\GL_N(\mathbf F_2)$ for $N\ge3$ is perfect, since
$1+E_{ij}=[1+E_{ik},1+E_{kj}]$ for distinct indices, and the
elementary matrices generate it.  It acts faithfully and doubly
transitively on nonzero vectors.  A nontrivial normal subgroup
is therefore transitive.  The stabilizer of a vector $v$ has
the abelian normal subgroup
$U_v=\{1+v\otimes\phi:\phi(v)=0\}$, and the conjugates of $U_v$
generate the whole group.  Modulo a transitive normal subgroup,
every conjugate of $U_v$ has the same image: move $v$ by an
element of the normal subgroup and then use its stabilizer.
The quotient is thus abelian, and perfectness makes it trivial.

Goursat's elementary common-quotient description now says that
the boundary image on $Q\oplus T_0$ contains
$\GL(Q)\times1$: a nontrivial common quotient would be the
simple group $\GL(Q)$, whose large prime cannot occur on $T_0$.
Let $B$ be the inverse image of $\GL(Q)\times1$.
It fixes $T_0$ pointwise and maps onto $\GL(Q)$.
Its kernel $R$ is a submodule of
$\Hom(Q,T_0)=Q^*\otimes T_0$, hence
$R=Q^*\otimes T'$ for a subspace $T'\subseteq T_0$.
Modulo $R$, the subgroup is a graph of a cocycle with values in
$Q^*\otimes(T_0/T')$.  The elementary vanishing
$H^1(\GL_{80}(\mathbf F_2),Q^*)=0$ from
Lemma~\ref{lem:char2-group-inputs} makes it conjugate to the
zero cocycle.  Its inverse image is then
$R\rtimes\GL(Q)$, and in particular contains the required
block-diagonal $\GL(Q)$.  A quotient surjection alone would
not have supplied this last subgroup; the cocycle step does.
\end{proof}

\begin{theorem}[Conditional exact uniform product groups]
\label{thm:product-full-group-if-hasse-irreducible}
Fix $d\ge3$ on the surface $E\times E$ above.
If the degree-$d$ Hasse divisor is geometrically irreducible,
then both geometric and arithmetic groups are
\[
 \begin{cases}
 \GL_{9d^2-1}(\mathbf F_2),&d\text{ odd},\\
 \mathbf F_2^{\,3(9d^2-4)}\rtimes\GL_{9d^2-4}(\mathbf F_2),
       &d\text{ even},
 \end{cases}
\]
where in the second line the group is the full pointwise
stabilizer of the three known independent covectors.
Theorem~\ref{thm:product-Hasse-irred} below proves the required
geometric irreducibility for every $d\ge4$, thereby completing
the entire fixed-surface full-image problem together with the
three certified initial degrees.
\end{theorem}

\begin{proof}
Put $q=0$ for odd $d$ and $q=3$ for even $d$; these are the
entire invariant-covector dimensions by
Theorem~\ref{thm:product-char2-all-parities-fixed-quotient}.
Let $K$ be the common kernel and $N=\dim K=9d^2-1-q$.
Proposition~\ref{prop:product-hasse-reduced} and the hypothesis
permit Lemma~\ref{lem:irreducible-hasse-quotients-trivial}.
Every nonzero proper invariant subspace therefore contains $K$.
It follows that $K$ is irreducible (or that $V=K$ is irreducible
when $q=0$), and that $V$ has no invariant vector: an invariant
line could not contain $K$, whose dimension is at least $80$.

The image on $K$ is irreducible and generated by transvections,
by Theorem~\ref{thm:product-all-degree-transvections}.
The natural $80$-dimensional block of
Lemma~\ref{lem:product-degree-three-block} lies in $K$, because
a nontrivial natural module has no equivariant map to a trivial
module.  On $K$ that subgroup acts naturally on this block and
trivially on its complement.
No nondegenerate alternating form on $K$ can be invariant under
it.  The natural module of $\GL_{80}(\mathbf F_2)$ has no
invariant bilinear form: elementary transvections first kill
every off-diagonal matrix entry, using a third index, and then
every diagonal entry.  Its pairing with a trivial module is
also zero.  Hence the entire natural block would lie in the
radical of any such form.

The alternatives in the opened
\cite[Theorem~3.6, p.~22]{KantorSeress2001} are therefore excluded,
and the image on $K$ is $\GL(K)$.
If $q=0$ this is the desired full group.
If $q=3$, the translation-kernel and $H^1$ argument in
Theorem~\ref{thm:product-char2-two-full-groups} applies verbatim.
The kernel is $K\otimes U$ for some $U\subseteq Q^*$, where
$Q=V/K$ is constant; vanishing of the natural-module first
cohomology conjugates the group to
$(K\otimes U)\rtimes\GL(K)$.
Its invariant-vector space is $U^\perp$.
The already proved vanishing $V^G=0$ forces $U=Q^*$, so all
translations occur.  The arithmetic image is between this
geometric group and the same pointwise stabilizer, because the
three covectors are defined over $\mathbf F_2$.
It is consequently the same group.
\end{proof}

\begin{remark}[Why reducedness alone was insufficient]
Generic stable loss one proves reducedness, but does not by
itself exclude two reduced components meeting in codimension
two. The subsequent proof of
Theorem~\ref{thm:product-Hasse-irred} supplies the missing
irreducibility by mixed tacnodal unions and smooth points
joining all restriction factors. It does not assume a generic
symmetric-matrix rank bound for the special Cartier matrices.
\end{remark}
